% 本文件由 LaTeX 排版 Agent 根据有效 translation.json 自动生成。
% author=Tertullian; work_id=0311; title=针对异端的抗辩书

\chapter{针对异端的抗辩书}

% structure: p0001 source=h1 target=omitted reason=page-title-already-rendered-by-parent

% structure: p0002 source=h2 target=section reason=relative-heading-rank; base=h2

\section{第一章 绪论。异端必然存在，甚至比比皆是；它们是对信德的考验。}

我们所处的时代特质，甚至促使我们提出这样的劝诫：我们不应因（充斥的）异端而惊愕，也不应因它们的存在而诧异，因为早已预言它们必将发生；亦不应因它们颠覆某些人的信德而困惑，因为它们的终极目的，正是通过给信德提供考验，也赋予它获得\Quote{赞许}的机会。\BibleRef{格前 11:19}因此，许多人因异端如此盛行而跌倒，这种见怪实属无据且轻率。倘若异端根本不存在，他们的跌倒该有多大呢？当一件事被注定必定发生时，它便获得了其存在的原因。这确保了它赖以存在的能力，以至于它不可能不存在。\par

% structure: p0004 source=h2 target=section reason=relative-heading-rank; base=h2

\section{第二章 热病与异端之间的类比。异端不足为奇：它们的力量源自人信德的软弱。它们并不拥有真理。以拳击手和角斗士的比喻说明。}

以热病这一相似情况为例，它在所有其他导致死亡和剧痛的（生命）结局中被指定了毁灭人的位置：我们既不因它的存在而惊讶（因为它就在那里），也不因它消耗人的生命而惊讶（因为这就是它存在的目的）。同样，对于异端，它们被产生出来是为了削弱和消灭信德，既然我们因它们拥有这种能力而感到恐惧，我们应先恐惧它们存在这一事实；因为只要它们存在，它们就有能力；只要它们有能力，它们就继续存在。但热病，由于其起因和力量皆为邪恶，尽人皆知，我们宁可厌恶它而不惊奇它，并尽力防范它，尽管我们无法根除它。然而，有些人却宁愿对异端感到惊奇，它们带来永死和更猛烈之火的炙热，只因它们具有这种能力，而不在有逃脱之机时避开它们的能力：但若（人们）不再惊异于异端拥有如此能力，它们便毫无能力。因为事情无非这样发生：要么人们在惊异时落入陷阱，要么由于被困而珍惜他们的惊讶，仿佛异端因其拥有的某种真理而如此强大。无疑，邪恶本身拥有任何力量本是一件奇妙的事，要不是异端在信德不坚的人身上强大。在拳击手和角斗士的打斗中，一般而言，一个人取胜不是因为他强壮，失败也不是因为他不够强壮，而是因为被击败的是一个毫无力量的人；事实上，这个得胜者后来若与真正强大的人交手，便会垂头丧气地退出比赛。完全同样，异端从个人的软弱中汲取它们的力量——每当遇到真正坚强的信德时，它们便毫无力量。\par

% structure: p0006 source=h2 target=section reason=relative-heading-rank; base=h2

\section{第三章 软弱的人轻易成为异端的猎物，异端从人类普遍软弱中汲取力量。杰出人物也曾从信德中跌倒：\Person{撒乌耳}、\Person{达味}、\Person{撒罗满}。基督的恒心。}

确实，性格较弱的人通常会被某些被异端所迷惑的人（在信心上）所建立，以至于自己也陷入毁灭。他们问：怎么会发生这样的事，那个在教会中最忠信、最明智、最受认可的女人或男人，竟然投向了另一边？提出这种问题的人，难道不是自己就已经回答了吗？即那些异端能够使之背道的人，本就不应被视为明智、忠信或受认可的人？这又是一件稀奇的事吗，即一个曾经受认可的人后来又退步了？\Person{撒乌耳}曾超越众人为善，后来却因嫉妒而败亡。\Person{达味}是个善人，\BibleQuote{合主心意}，\BibleRef{撒上 13:14}后来却犯下杀人和奸淫之罪。\EditorialNote{撒慕尔纪下 11章}\Person{撒罗满}蒙主赐予一切恩宠和智慧，却被女人引向偶像崇拜。\BibleRef{列上 11:4}因为只有天主子才能无玷地坚持到底。但倘若一位主教、一位执事、一位寡妇、一位贞女、一位教师、甚或一位殉道者，从（信仰）规则上跌倒，难道异端就会因此显得拥有真理吗？我们是以人来证明信仰，还是以信仰来证明人？无人是明智的，无人是忠信的，无人是尊贵的，除非是基督徒；也无人是基督徒，除非他坚持到底。\BibleRef{玛 10:22}你，作为一个人，从外表认识别人。你按你所见的去想。而你只能看到眼睛所及之处。但（圣经）却说：\BibleQuote{主的眼目高远。}\BibleRef{耶 32:19}\BibleQuote{人看外表，但主看人心。}\BibleRef{撒上 16:7}\BibleQuote{主（注视并）认识属于祂的人；}\BibleRef{弟后 2:19}并且\BibleQuote{那不（由我天父）栽种的植物，祂要连根拔起；}\BibleRef{玛 15:13}并且\BibleQuote{首先的，}如祂所显示，\BibleQuote{将成为最后的；}\BibleRef{玛 20:16}并且祂拿着\BibleQuote{簸箕在手，要扬净祂的禾场。}\BibleRef{玛 3:12}任那轻浮信德的糠秕在每一次诱惑的风中飘散，那将被收在主仓中的麦粒就会更加纯净。难道不是有一些门徒因跌倒而离开了主自己吗？\BibleRef{若 6:66}然而其余的人并没有因此认为他们也必须远离跟随祂，而是因为他们知道祂是生命之言，是从天主来的，他们便一直与祂同在，直到最后，在祂柔和地问他们是否也要离去之后。\BibleRef{若 6:67}像\Person{菲革路}、\Person{赫摩革乃}、\Person{斐理托}和\Person{依默纳约}这样的人离弃了祂的宗徒，这还算小事；基督的出卖者自己就是宗徒之一。我们惊讶于祂的教会被一些人离弃，然而我们效法基督自己所受的苦，正表明我们是基督徒。\Person{圣若望}说：\BibleQuote{他们从我们中间出去，却不属于我们；如果他们属于我们，他们必会留在我们中间。}\par

% structure: p0008 source=h2 target=section reason=relative-heading-rank; base=h2

\section{新约中关于异端的警告。引用多处经文。这些经文暗示了陷入异端的可能性。}

但我们更当铭记主的话语和宗徒的书信；因为他们都已预先告诉我们将有异端，并事先警告我们要躲避它们；既然我们不因它们存在而惊慌，也就同样不应惊讶它们有能力做出那使得我们必须躲避它们的事。主教导我们，有许多\BibleQuote{披着羊皮的凶狼}要来。\BibleRef{玛 7:15}那么，这些羊皮是什么，岂非基督徒职业的外表？凶狼是谁，岂非那些潜伏在内、要毁灭基督羊群的欺骗性的感官和灵？假先知是谁，岂非欺诈的未来预言者？假宗徒是谁，岂非传播假冒福音的宣讲者？还有那敌基督者，无论现在还是将来，岂非那些悖逆基督的人？现今的异端将以其教义的歪曲撕裂教会，其程度不亚于那日敌基督以残忍攻击逼迫教会，只是逼迫使七人成为殉道者，而异端只造就背教者。因此\BibleQuote{必须要有异端，好使那些受认可的人被显露出来，}\BibleRef{格前 11:19}即在逼迫中坚持到底的人，以及未迷失道路而进入异端的人。因为\Emph{宗徒}的意思并非指那些将信经换成异端的人应被视为\Emph{受认可的}；尽管他们曲解他的话以利于自己，当他在另一处说：\Quote{查验一切，持守那善好的；}好像一个人经过一切错误的查验后，不会因错误而固执地选择某种邪恶之事似的。\par

% structure: p0010 source=h2 target=section reason=relative-heading-rank; base=h2

\section{异端与分裂、纷争一样，为\Person{圣保禄}所不悦；他论及异端的必要性，并非作为善事，而是出于天主的旨意，为训练和证明基督徒信德的有益考验。}

此外，当他责备那些无疑是恶事的纷争与分裂时，他立即也加上了\Emph{异端}。他既将异端列于恶事之列，当然也承认它是恶的；而且，它确实是更大的恶，因为他告诉我们，他相信他们的分裂与纷争，是基于他得知\BibleQuote{也必须要有异端} \BibleRef{格前 11:19}。他向我们表明，正是由于预见到更大的恶，他才轻易相信了较轻恶事的存有；而且，他远未相信（在诸如此类的恶事中）异端是善的，他的目的是要预先警告我们，不要对更恶劣的试探感到惊讶，因为（他说）它们\BibleQuote{要显明那些是经得起考验的} \BibleRef{格前 11:18}，换言之，即是那些不能被他们腐化的人。简言之，既然整段经文旨在维护合一、制止分裂，而异端既不亚于分裂与纷争，同样使人脱离合一，无疑他将异端与分裂和纷争归于同一谴责之下。如此，他便使那些陷于异端的人成为\Quote{不经得起考验的}；更因他以斥责劝勉人远离他们，教导他们应当\BibleQuote{都说一样的话，有一样的心思} \BibleRef{格前 1:10}，这正是异端所不允许的事。\par

% structure: p0012 source=h2 target=section reason=relative-heading-rank; base=h2

\section{第六章 异端者自我定罪。异端是自我意志，而信仰是意志服从天主的权威。\Person{阿佩莱斯}之异端。}

然而，在此点上我们不再赘述，因为正是同一位保禄在\WorkTitle{致迦拉达人书}中将\BibleQuote{异端}列于\BibleQuote{肉性的作为}之中 \BibleRef{迦 5:20}；他也提醒弟铎说，\BibleQuote{异端者，在初次劝诫之后，必须予以弃绝}，理由是\BibleQuote{这样的人已背道而驰，且犯罪，自己定自己的罪} \BibleRef{铎 3:10}。事实上，在几乎每一封书信中，当他命令我们躲避假道理时，他都严厉谴责异端。这些异端的实际效果是假道理，希腊文称为\OriginalTerm{异端}{\GreekText{αἵρεσις}}，这个词的意义是指一个人在自己教导（别人）或自己接受（它们）时所做出的那种\Emph{选择}。正因为如此，他称异端者为\Emph{自我定罪}，因为他自己选择了那使他被定罪的东西。然而，我们不可凭自己的意志执著任何对象，也不可挑选别人凭私意引入的东西。我们拥有主宗徒的权威；因为他们也不是凭自己选择引入任何东西，而是忠实地将自基督所受的道理传给了（人类的）万邦。因此，即使\BibleQuote{从天上来的天使传别的福音}（不同于他们的），他也必为我们所咒诅。圣神甚至已预见到，在某一位（名叫）\Person{菲卢墨涅}的童贞女那里将有一个欺骗的天使，\BibleQuote{装作光明的天使} \BibleRef{格后 11:14}，\Person{阿佩莱斯}被她的奇迹和幻象所引导，（当他）引入他的新异端时。\par

% structure: p0014 source=h2 target=section reason=relative-heading-rank; base=h2

\section{第七章 外教哲学是异端之根源。偏离基督信仰与旧时外教哲学体系间的关联。}

这些都是\Quote{人的学说}和\Quote{魔鬼的教训}\BibleRef{弟前 4:1}，是为了这世界智慧之灵的发痒的耳朵而产生的；主称这些为\Quote{愚妄}，并且\Quote{拣选了世上愚妄的事}来使哲学本身也感到羞愧。因为哲学正是世界智慧的材料，是鲁莽解释天主的本性与安排者。事实上，异端本身就是由哲学煽动的。从这个源头产生了\OriginalTerm{永世}{AEons}，以及我不知道的无数形式，还有属于\Person{瓦伦提努}系统中的\Quote{人的三重性}——而瓦伦提努是\Person{柏拉图}学派的人。从同一源头产生了\Person{马西昂}的更好的神，带着他的一切宁静；他来自\OriginalTerm{斯多葛派}{Stoics}。此外，灵魂会死的观点为\OriginalTerm{伊壁鸠鲁派}{Epicureans}所持守；而否认身体复活的则来自所有哲学家的综合学派；再者，当物质被等同于天主时，你便有了\Person{芝诺}的教导；当有论及火神的任何学说时，\Person{赫拉克利特}便登场了。相同的问题被异端者和哲学家们一论再论；相同的论点被涉及。邪恶从何而来？为何允许它存在？人的起源是什么？他是如何而来的？此外还有\Person{瓦伦提努}最近提出的问题——天主从何而来？他用这个答案来解决：来自\OriginalTerm{思想}{enthymesis}和\OriginalTerm{流产}{ectroma}。可悲的\Person{亚里士多德}！他为这些人发明了辩证法，一种搭建与拆毁的艺术；这门艺术在其命题上如此闪烁其词，在其推测上如此牵强，在其论证上如此严苛，如此富于争辩——甚至连它自己也感到困惑，收回一切，实则讨论无物！那些\Quote{寓言与无尽的族谱}\BibleRef{弟前 1:4}、\Quote{无益的问题}\BibleRef{铎 3:9}以及\Quote{像毒癌一样扩散的话语}\BibleRef{弟后 2:17}是从何处滋生出来的？从这一切中，当宗徒想要约束我们时，他明确地指出了\Emph{哲学}，作为我们应该警惕的事物。他写信给\Place{哥罗森}人时说：\Quote{你们要谨慎，恐怕有人用哲学和虚妄的欺骗，依照人的传统，相反圣神的智慧，诱惑你们。}他曾到过\Place{雅典}，并在与哲学家的谈话中了解了那种自称认识真理却只败坏真理的人性智慧，并且它自身由于各种互相抵触的派别而分裂为自身多样的异端。\Place{雅典}与\Place{耶路撒冷}究竟有何相干？\Place{学园}与教会之间有何和谐？异端者与基督徒之间有何共同之处？我们的教导来自\Quote{\Person{撒罗满}的门廊}，他亲自教导说：\Quote{当以心地纯朴寻求上主。}\BibleRef{智 1:1} 摒弃所有制造出\OriginalTerm{斯多葛派}{Stoic}、\OriginalTerm{柏拉图派}{Platonic}和辩证法的混杂基督教的企图！拥有耶稣基督之后，我们不再需要好奇的争辩；享受福音之后，不再需要探究！有了我们的信德，我们不再渴望进一步的信仰。因为这是我们的主要信德：除此外我们不应相信任何事物。\par

% structure: p0016 source=h2 target=section reason=relative-heading-rank; base=h2

\section{第八章 基督的话\Quote{寻求，就必寻见}并不为偏离信仰的异端提供根据。基督对犹太人所说的一切话都是为了我们，虽然不是作为具体命令，而是作为要应用的原则。}

我现在来到这一点，这一点（被我们的弟兄和异端者都主张）。我们的弟兄把\Emph{它}作为一个借口来从事好奇的探究，而异端者则坚持用它来引入（他们不信的）苛刻。他们引用说：\Quote{寻求，就必寻见。}\BibleRef{玛 7:7} 让我们记住主说这话的时间。我想是在祂教导的最初，那时所有人仍然怀疑祂是否是基督，甚至\Person{伯多禄}还没有宣认祂是天主子，而\Person{若翰}（洗者）实际上已经对祂失去了确信。因此，当时说\Quote{寻求，就必寻见}是很有道理的，因为那时还需要对那尚未被认识的主进行探究。此外，这话\Emph{曾说过}是就犹太人说的。因为整个责备的事正是针对他们的，因为他们有（启示）可以在那里寻求基督。\par

祂说：\Quote{他们有梅瑟和厄里亚，} \BibleRef{路 16:29} ——即法律和先知，他们宣讲基督；又如祂在另一处清楚地说：\Quote{你们查考圣经，因你们想在其中得永生；正是这些为我作证；} \BibleRef{若 5:39} 这也正是\Quote{寻找，你们就会找到}的意思。因为显然接下来的话也是对犹太人说的：\Quote{敲门，就给你们开门。} \BibleRef{玛 7:7} 犹太人从前曾与天主立约；但后来因他们的罪被遗弃，便开始没有天主。外邦人则相反，从未与天主立约；他们不过是\Quote{桶里的一滴}和\Quote{场上的灰尘} \BibleRef{依 40:15}，始终在门外。现在，那从未在门内的人，怎能敲他从不在的地方的门呢？他既没有通过入门或驱逐而经过任何门，又怎会知道哪扇门呢？难道不是那意识到自己曾住在里面又被赶出来的人，才（可能）找到门并敲门吗？同样地，\Quote{求，就会得到} \BibleRef{玛 7:7} 是对那知道应向谁求、也得了谁应许的人说的，也就是说：\Quote{亚巴郎、依撒格、雅各伯的天主。} 而外邦人既不认识祂，也不知道祂的任何应许。因此祂是对以色列说的：\Quote{我被派遣，只是为了以色列家迷失的羊。} \BibleRef{玛 15:24} 祂还没有\Quote{把儿女的饼扔给狗}；还没有吩咐他们\Quote{不要走外邦人的路。} \BibleRef{玛 10:5} 直到最后，祂才教导他们\Quote{去使万民成为门徒，因父及子及圣神之名给他们授洗，} \BibleRef{玛 28:19} 那时他们即将领受\Quote{圣神，那位护慰者，祂要引领你们进入一切真理。} \BibleRef{若 16:13} 这件事也有助于同样的结论。如果被立为外邦人的老师的宗徒们，自己还要以护慰者为师，那么对我们说\Quote{寻找，你们就会找到}就更没有必要了——因为不需研究，我们的教导将由宗徒们而来，而宗徒们则由圣神而来。的确，主的一切话都是向众人宣讲的，经过犹太人的耳朵传到我们这里。然而，其中大部分是对\Emph{犹太人}说的，因此它们并不专为我们而设的教导，而是\Emph{更}像是一个榜样。\par

% structure: p0019 source=h2 target=section reason=relative-heading-rank; base=h2

\section{第九章 寻找确定真理的训诫。当我们发现这真理时，应当满足。}

我现在故意放弃这一论证立场。就算承认\Quote{寻找，你们就会找到}这句话是（平等地）对所有人说的。即便如此，人仍需谨慎地按照（那）统辖一切诠释的原则来确定这些词句的意义。（然而）没有一句神性之言是那样孤立和散漫，以致只需坚持其\Emph{字词}，而可忽略其\Emph{上下文联系}。但我从一开始就提出（这个主张）：基督所教导的某件事，且是确定的一件事，是外邦人务必相信的，并且为此目的要去\Quote{寻找}，以便他们当\Quote{找到}时能够相信。然而，对于那被教导为唯一确定之事，不能有不确定的寻找。你必须\Quote{寻找}直到你\Quote{找到}，找到后就要相信；除了持守所信的，你无需再做别的事，只要你此外还相信：在找到并相信了那命令你不可寻求祂所教导之外的事的那一位所教导的事之后，再无其他可相信，因此也无其他可寻求的。确实，当有人对此怀疑时，我们将提供证据，证明我们拥有那由基督所教导的道理。与此同时，我凭我们的证据如此自信，以致我预先提醒某些人：不要\Quote{寻求}超过他们所信的——这正是他们本应寻求的事，即如何避免不顾理性规则地解释\Quote{寻找，你们就会找到}。\par

% structure: p0021 source=h2 target=section reason=relative-heading-rank; base=h2

\section{第十章 当一个人相信时，他已成功地找到了确定真理。异端者总是提出许多事物供人空论，但我们不应总是寻求。}

如今，这句格言的理由包含在三个方面：在质料方面，在时间方面，在界限方面。在质料方面，你必当考虑你在寻求\Emph{什么}；在时间方面，\Emph{何时}寻求；在界限方面，\Emph{寻求多久}。那么，你所应当\Quote{寻求}的，乃是基督所教导的，（你必当继续寻求）当然直到你寻不着为止——直到你确实寻见为止。但你相信时，你已经寻见了。因为你若未曾寻见，就不会相信；同样，除非为了寻见，你也不会寻求。因此，你寻求的目的是为了寻见；而寻见目的乃是为了相信。你已通过相信，阻止了进一步寻求与寻见的拖延。你寻求的果实已为你确定了这一界限。为你设定这界限的正是那一位，祂不愿意你相信任何祂所教导之外的事物，因此也不愿你去寻求它。然而，正因为有许多人各自教导了不同的事，我们便有义务继续寻求，只要我们能寻见什么；那么（依此逻辑）我们将永远寻求，而从未相信任何事物。因为寻求的终点在哪里？相信的止境在哪里？寻见的完成在哪里？（要与）\Person{马尔西翁}一起吗？但\Person{瓦伦廷}也向我们提出这格言：\Quote{寻求，就必寻见。}（那么要与）\Person{瓦伦廷}一起吗？看，\Person{阿佩莱斯}也要用同样的引文攻击我；\Person{厄比翁}、\Person{西满}以及所有其他人，也都没有别的论据来引诱我、拉拢我。这样，我将无处立足，仍不断遭遇那句挑战：\Quote{寻求，就必寻见，}仿佛我没有安息之所；仿佛我从未寻见基督所教导的——那应当寻求、必须相信的。\par

% structure: p0023 source=h2 target=section reason=relative-heading-rank; base=h2

\section{第十一章 相信之后，寻求应当停止；否则必然导致对所信之事的否认。我们的信仰不以其他事物为目标。}

若没有过失，错误便可免罚；尽管错误本身也是过失。我说，当一个人（故意）不放弃任何事物时，他漫游是免罚的。然而，若我相信了我所应当相信的，后来又以为有新的东西要寻求，我自然期望还有别的可寻见，尽管我决不应存此期望，除非因为我要么并未相信——尽管我表面上是信徒——要么已停止相信。我若如此背弃信仰，便被证明是信仰的否认者。我一次且永远地说：没有人寻求，除非他要么从未拥有，要么失去了（他所寻求的）。（福音中的）老妇人丢失了十块银钱中的一块，所以她寻找它；\BibleRef{路 15:8}但当她寻见时，便不再寻找。邻居没有饼，所以他敲门；但门一开，他得到了饼，便不再敲了。\BibleRef{路 11:5}寡妇不断恳求法官听她申诉，因为她未获准；但她的案件得到审理后，她便不再作声。\BibleRef{路 18:2}因此，无论寻求、敲门还是祈求，都有限度。祂说：\Quote{凡祈求的，}祂说，\Quote{就必得着，叩门的，就必给他开门，寻求的，就必寻见。}\BibleRef{路 11:9}愿那永远寻求却从未寻见的人被除去！因为他寻求的地方什么也寻不着。愿那永远敲门却总不被打开的人被除去！因为他敲的门后面无人（开启）。愿那永远祈求却总不被垂听的人被除去！因为他向不听的人祈求。\par

% structure: p0025 source=h2 target=section reason=relative-heading-rank; base=h2

\section{第十二章 对神圣知识的正当寻求，永不会不合时宜或过度，且总在信仰的规范之内。}

至于我们，虽然仍需寻求，而且\Emph{始终}如此，但我们的寻求应在何处进行？在异端者中间吗？那里的一切都与我们的真理相异、相敌，我们也被禁止接近他们。有哪个奴隶从陌生人，更不用说从主人的仇敌那里寻找食物？有哪个士兵期望从盟国国王——我几乎可以说不友好的国王——那里得到赏赐和薪饷？除非他是一名逃兵、一名叛逃者和一名反叛者。就连那老妇人也只在自己的家中寻找银钱。那坚持不懈敲门的人也只在邻居的门前敲打。寡妇也不是向仇视的法官（尽管严苛）上诉。没有人从导向毁灭的事物中获得教导。没有人从全然黑暗的地方得到光照。因此，让我们的\Quote{寻求}停留在属于我们的事物中，来自属于我们的人那里；并针对属于我们的事物——唯独那些可以在不损害信仰规范的情况下成为探究对象的事物。\par

% structure: p0027 source=h2 target=section reason=relative-heading-rank; base=h2

\section{第十三章 信经或信仰规范的摘要。信徒从未对信经提出疑问。异端者鼓励并延续脱离基督教导的独立思想。}

现在，关于这\OriginalTerm{信德规则}{rule of faith}——好让我们从此点承认我们所辩护的是什么——它，你得知道，就是那规定信仰：只有一个天主，而祂不是别的，就是世界的创造者，祂通过自己的圣言——首先发出的——从无中产生了万物；这圣言被称为祂的圣子，\Emph{并且}，在天主的名下，以\Quote{多重方式}被先祖们看见，在先知们中时常被听到，最后被圣父的神与德能带下进入童贞玛利亚，在她的子宫中成为血肉，并由她出生，出来作为耶稣基督；从此祂宣讲新法律和天国的新许诺，行奇迹；被钉十字架后，第三天复活；然后升天，坐在圣父的右边；派遣了圣神的德能代替自己来领导那些相信的人；将带着光荣来临，接圣人到永生和天堂许诺的享受中，并惩罚恶人到永火，在两种人都复活之后，连同他们的肉身的恢复。这规则，正如将要证明的，是由基督教导的，且在我们中间不引起任何问题，除了异端所引进的、以及使人成为异端者的问题。\par

% structure: p0029 source=h2 target=section reason=relative-heading-rank; base=h2

\section{第十四章。好奇心不该越过信德规则。不安的好奇心，异端的特征。}

然而，只要它的形式以适当次序存在，你可以随意寻求和讨论，并按你的心意放纵好奇心，在任何你觉得悬疑或晦暗不明的事情上。你手边无疑有一些博学的弟兄，蒙受知识的恩宠，一些经验丰富的人，一些你亲近的、和你一样好奇的熟人；虽然他自己和您一样是个寻求者，但毕竟他会很清楚，你最好保持在无知中，免得你得知你不该知道的事，因为你已经获得了你应当知道的知识。\Quote{你的信德，}他说，\Quote{救了你}\BibleRef{路 18:42}，而不是\Emph{钻研}你在圣经上的技艺。如今，信德被存放在规则中；它有法律，并且在遵守中得有救恩。然而，技艺在于好奇的巧术，其光荣只是来自窍门的敏捷。让这样的好奇技艺让位给信德；让这样的光荣屈从救恩。无论如何，让他们要么放弃喧闹，要么安静。在（信德）规则之外一无所知，就是知道一切。假设异端不是真理的仇敌，以至于我们没有被预先警告要躲避他们，那么与那些自称仍在寻求的人达成一致，将是怎样的行为？因为他们若仍在寻求，就尚未找到任何确定的事；因此，无论他们暂时似乎持有什么，他们在继续寻求时，就暴露了自己的怀疑论。因此，你既然照他们的方式寻求，效仿那些自己永远在寻求的人，一个怀疑者对怀疑者，一个动摇者对动摇者，必然会被\BibleQuote{若是盲人领盲人，两人必掉进坑里}\BibleRef{玛 15:14}。但是，当他们为了欺骗我们，假装仍在寻求，以便通过焦虑同情的暗示，把他们的论文强塞给我们——简言之，当他们（在获得进入我们之中后）立刻坚持我们必须探究他们惯于提出的那些问题，那么，我们在道德义务上就应当拒绝他们，好让他们知道，我们所否认的不是基督，而是他们自己。因为他们既然是寻求者，就尚未有固定的主张；既然没有固定主张，就尚未相信；既然尚未相信，就不是基督徒。但即使他们有自己的主张和信仰，他们仍然说为了讨论必须探究。然而，在讨论之前，只要他们还把它当作探究的对象，他们就否认自己尚未相信所承认的事。因此，当人们甚至连他们自己也承认不是基督徒时，对我们来说，他们更不是基督徒！当他们用谎言向我们推荐真理时，他们所赞助的是什么样的真理？好吧，但他们实际上是在谈论圣经，并从圣经中推荐他们的意见！他们确实如此。除了从信仰的记录中，他们还能从哪里获得关于信仰之事的论据呢？\par

% structure: p0031 source=h2 target=section reason=relative-heading-rank; base=h2

\section{第十五章。异端者不准许从圣经中辩论。事实上，圣经不属于他们。}

因此我们来到了我们立场（的要点）；因为我们瞄准了这一点，并且我们在我们讲话的引言（刚刚完成）中为此作了准备——好让我们现在可以对我们的对手向我们挑战的争论进行交锋。他们提出圣经，并凭借他们的这种无礼，他们立刻影响了一些人。然而，在交锋本身中，他们使强壮者疲惫，捕获软弱者，并以怀疑打发摇摆者。因此，我们反对他们采取这一比其他一切都优先的步骤，即不让他们参与任何关于圣经的讨论。\par

如果他们的资源在于这些（圣经），在他们能使用它们之前，应当清楚地看出圣经的所有权属于谁，以便根本没有资格享受这一特权的人不得被允许使用它们。\par

% structure: p0034 source=h2 target=section reason=relative-heading-rank; base=h2

\section{第十六章。宗徒对将异端者排除在圣经使用之外的认可。按照宗徒的话，异端者不应被辩论，而应被劝诫。}

或许有人会认为我提出这一立场是为了补救我案件中的不信任，或是出于以其他方式进入争战的愿望，但若不是我有支持的理由，尤其是这一点：我们的信德应服从宗徒，他禁止我们参与\Quote{问题}，或去倾听新奇的说法，\BibleRef{弟前 6:3} 或与异端者\Quote{在第一次和第二次劝诫之后}来往，\BibleRef{铎 3:10} ——请注意，不是讨论之后。他用这种方式禁止讨论，指定与异端者打交道的目的乃是\Emph{劝诫}，而且是\Emph{第一次}劝诫，因为他不是基督徒；这样，他就不至于像一个基督徒那样，似乎需要一而再、再而三的纠正，并且\Quote{在两三个证人面前}，\BibleRef{玛 18:16} 因为他应当被纠正，正是基于他不应当被争论；其次，因为关于圣经的争论显然只能产生一种效果，即扰乱人的胃或头脑。\par

% structure: p0036 source=h2 target=section reason=relative-heading-rank; base=h2

\section{第十七章　异端者事实上并不使用圣经，而只是滥用圣经。他们与你之间没有共同基础。}

你们的这个异端并不接受某些圣经；而它所接受的，它也为了达到自己的目的，通过增减加以歪曲；它所接受的，也不是完整地接受；但即使它在一定范围内完整地接受了某些书卷，它也仍然通过多种解释的诡计加以歪曲。真理既因对其意义的窜改而受到反对，也因其文本的败坏而受到反对。他们虚妄的僭越必须拒绝承认那些他们被驳斥的（著作）。他们依靠那些他们错误拼凑的、并因含混而选择的著作。尽管在圣经方面极为熟练，你将毫无进展，因为你所坚持的一切都被对方否认，而你所否认的一切都被（他们）坚持。至于你自己，你确实只会浪费口水，并从他们的亵渎中得到烦恼。\par

% structure: p0038 source=h2 target=section reason=relative-heading-rank; base=h2

\section{第十八章　从圣经出发的任何讨论，给信德软弱的人带来巨大祸害。异端者绝不会通过这种过程得到说服。}

但至于你为了他而进入圣经讨论的那个人，目的是当他被怀疑困扰时坚固他，（请问）他会倾向真理，还是倾向异端见解呢？正是因为他看到你没有取得任何进展，而对方在否定和辩护方面与你旗鼓相当（或者至少处于类似地位），他将因讨论而在不确定中离开，不知道哪一方应被判定为异端。因为，毫无疑问，他们也能把这些事反驳到我们身上。这确实是一个必然的结果：他们甚至会说，圣经的窜改和错误诠释是我们自己引入的，因为他们和我们一样坚持真理在他们一边。\par

% structure: p0040 source=h2 target=section reason=relative-heading-rank; base=h2

\section{第十九章　在讨论异端时，上诉不在于圣经。圣经只属于那些拥有信德准则的人。}

因此，我们的上诉不可指向圣经；也不得接受在那些胜利要么不可能、要么不确定、要么不够确定的点上争论。但是，即使从圣经出发的讨论不会导致双方旗鼓相当，（然而）事物的自然顺序会要求首先提出这一点，这是我们现在唯一必须讨论的：\Quote{那圣经所属的信德本身，在谁那里？从何处、借着谁、在何时、传给谁，这使人成为基督徒的准则被传授下来？}因为无论在哪里，只要显明真正的基督徒准则和信德所在，那里也就有真正的圣经及其解释，以及所有基督徒的传承。\par

% structure: p0042 source=h2 target=section reason=relative-heading-rank; base=h2

\section{第二十章　基督首先传授了信德。宗徒们传播了它；他们建立了教会作为它的宝库。因此，那信德是宗徒的，它从宗徒们，通过宗徒的教会传承下来。}

我们的主基督耶稣（愿祂容忍我一时如此表达！），无论祂是谁，无论祂是哪位天主的圣子，无论祂是人与天主由何种质料构成，无论祂是何种信仰的导师，无论祂是何等赏报的许诺者，当祂在世时，亲自宣告了祂是什么、祂曾是什么、祂所执行的圣父的旨意是什么、祂所规定的人的本分是什么；（这一宣告，）或是公开对民众，或是私下对祂的门徒，祂从门徒中拣选了十二位为首者常伴左右，\BibleRef{谷 4:34}并预定他们成为万民的导师。因此，在其中一个被除去之后，祂在回归圣父时命令其余十一人：\BibleQuote{你们要去使\Emph{所有}民族成为门徒，因圣父、圣子及圣神之名给他们授洗。}\BibleRef{玛 28:19}于是宗徒们立即遵行，这一称号正指出他们是\Quote{\Emph{被派遣者}}。他们依据大卫圣咏中预言的权威，以抽签方式选了玛弟亚为第十二位，以取代茹达斯；接着，他们获得了圣神所应许的德能，行奇迹并宣讲；在犹大境内首先为对耶稣基督的信仰作证，并在各处建立教会后，他们便出发走向世界，向万民宣讲同一信仰的同一道理。他们同样在各城建立了教会，所有其他教会相继从这些教会领受了信仰的传统和道理的种子，并每日领受，以便成为教会。事实上，正是因此，它们才能自视为宗徒的，作为宗徒教会的后裔。每一事物必然追溯其起源以归类。因此，教会虽众多且伟大，却只包含由宗徒建立的那唯一的原始教会，所有教会皆由此而生。如此，所有教会都是原始的，都是宗徒的，同时它们因和平的共融、手足之情和款待的纽带而被证明是一体的——这些特权除了同一奥迹的同一传统之外，别无其他规则所指引。\newline{}\par

% structure: p0044 source=h2 target=section reason=relative-heading-rank; base=h2

\section{第二十一章 凡通过教会从宗徒而来的全部道理都是真理，宗徒由天主藉基督教导；凡无如此神圣起源和宗徒传统可循的意见，皆\Emph{当然}为假。}

由此，我们制定我们的规则。既然主耶稣基督派遣了宗徒去宣讲，（我们的规则是）除了基督所委派的人以外，不应接受任何其他宣讲者；因为\BibleQuote{除了圣子，没有人认识圣父；除了圣子愿意启示的人，也没有人认识圣父。}\BibleRef{玛 11:27}圣子似乎也没有将圣父启示给宗徒以外的任何人，祂派遣宗徒去宣讲——当然，就是祂启示给他们的内容。现在，他们所宣讲的——换言之，基督所启示给他们的——正如我在此必须同样规定的那样，除了由宗徒亲自建立的这些教会本身，没有其他方式可以适当证明；宗徒们曾直接向他们宣布福音，既\Emph{亲口}（正如俗语所说），随后又通过他们的书信。因此，既然如此，显而易见，所有与宗徒教会——这些信仰的模子和原始源泉——相符的道理，必须被视为真理，因为它无疑包含着（所说的）教会从宗徒、宗徒从基督、基督从天主所领受的。反之，所有与基督和天主的教会及宗徒的真理相悖的道理，必须预先断定为虚假。那么，我们现在要证明：我们这已给出规则的道理，是否源于宗徒的传统，而所有其他\Emph{道理}是否\Emph{当然}出自虚假。我们与宗徒教会共融，因为我们的道理与\Emph{他们的}毫无差异。这就是\Emph{我们的}真理见证。\newline{}\par

% structure: p0046 source=h2 target=section reason=relative-heading-rank; base=h2

\section{第二十二章 驳斥试图推翻此信仰规则的做法。宗徒是真理的可靠传递者。他们起初已受充分教导，并在传递中保持忠信。}

然而，既然证据近在咫尺，若立即提出，便无事可处理，那么让我们暂时向对方让步，如果他们以为能找到某种否定这一规则的方法，仿佛我们拿不出任何证据似的。他们通常告诉我们，宗徒们并不知晓一切：（但在此）他们被同样的疯狂所驱使，由此转向完全相反的观点，宣称宗徒们确实知晓一切，却没有把一切交付给所有人——无论哪种情况，都将基督暴露于责备之下，因为祂派遣了要么过于无知、要么过于简单的宗徒。那么，一个心智健全的人，怎么可能认为那些由主指定为师傅（或导师）的人，在陪伴、做门徒、与主共处时，主使他们与自己密不可分，并且\Quote{当只有他们单独在一起时，祂常给他们解释}一切隐晦之事\BibleRef{谷 4:34}，告诉他们\Quote{那些奥秘被赐予他们去认识}\BibleRef{玛 13:11}，而这些奥秘是不允许百姓理解的？有什么是向\Person{伯多禄}隐藏的？他被称作\Quote{教会应建立在其上的磐石}，他也获得了\Quote{天国的钥匙}，并拥有\Quote{在天上和地上束缚和释放}的权柄？又有何物向\Person{若望}隐瞒？他是主最钟爱的门徒，曾靠在主的胸膛上，\BibleRef{若 21:20}主唯独向他指出\Person{犹达斯}是叛徒，并将祂的母亲\Person{玛利亚}托付给他，如同托付给自己的儿子一样？\BibleRef{若 19:26}难道主会打算叫那些人对什么无知吗？主甚至向他们显露了自己的光荣，与\Person{梅瑟}和\Person{厄里亚}一起，并且还有圣父的声音从天上传来？\BibleRef{玛 17:1}这并不是说主因此不赞成其余的人，而是因为\Quote{凭三个见证人的口，每一句话都应成立。}同样地，（我想，）那些人在复活之后，当主与他们同行时，蒙恩\Quote{得以解释一切圣经}，难道他们会无知吗？\BibleRef{路 24:27}无疑，主曾说过：\Quote{我还有许多话要告诉你们，但你们现在不能承受。}但即使那时祂也补充说：\Quote{当那真理之神来到时，祂要把你们引入一切真理。}\BibleRef{若 16:12}祂（由此）表明，那些祂应许借着真理之神获得将来对全部真理之通达的人，是没有一样事物是其所不知的。而且，祂确实履行了祂的应许，因为在\WorkTitle{宗徒大事录}中证实了圣神确实降临。现在，那些拒绝这圣经的人，既不属乎圣神，因为他们不能承认圣神已经赐予了门徒，也不能自称为教会，因为他们完全无法证明这个团体是何时、以何种裹布建立的。他们故意不为他们所主张的事提供任何证据，以免随着这些证据，那些他们虚伪杜撰的东西也带来有害的揭露。这于他们是如此重要。\par

% structure: p0048 source=h2 target=section reason=relative-heading-rank; base=h2

\section{第二十三章。宗徒并非无知。异端者关于圣伯多禄因受圣保禄责斥而不完美的妄论。圣伯多禄并非因教导错误而受责斥。}

现在，为了给宗徒们打上无知的印记，他们举出\Person{伯多禄}及其同伴受\Person{保禄}责斥的例子。他们说：\Quote{因此，}他们说，\Quote{他们有所欠缺。}（他们如此声称，）为的是由此建立他们的另一个论调，即后来可能有更圆满的知识临到（宗徒们），如同当\Person{保禄}责斥比他先来的人时他所分得的。我可以在此对那些拒绝\WorkTitle{宗徒大事录}的人说：\Quote{首先，你们必须向我们指明这位\Person{保禄}是谁——他成为宗徒之前是什么，以及他如何成为宗徒，}——他们在其他问题上对他是如此地大量使用。确实，他自己告诉我们，他在成为宗徒之前是个迫害者，\BibleRef{迦 1:13}但这对于任何一个在相信之前先考察的人来说还是不够的，因为连主自己也不为自己作证。\BibleRef{若 5:31}但如果他们的目的是要相信与圣经相反的事，那就让他们不信圣经吧。不过，他们仍然应该从他们所声称的\Person{伯多禄}受\Person{保禄}责斥这一事实中证明，\Person{保禄}加添了另一种形式的福音，与\Person{伯多禄}及其余的人先前所传的不同。但事实是，\Person{保禄}从迫害者转变为传道者，由弟兄们作为一位弟兄引荐给弟兄们——实际上，是由那些从宗徒们手中接受了信德的人引荐的。后来，正如他自己所叙述的，他\Quote{上耶路撒冷去见\Person{伯多禄}}\BibleRef{迦 1:18}，无疑是为了他的职务，并且基于共同信仰和宣讲的权利。现在，如果他的宣讲是相反的东西，他们当然不会对他从迫害者变为传道者感到惊讶；而且，他们也不会\Quote{光荣主}\BibleRef{迦 1:24}，因为\Person{保禄}显得是与主作对的人。他们甚至给了他\Quote{右手相交之礼}\BibleRef{迦 2:9}，作为与他一致的标记，并在他们之间安排了职务的分工，而不是福音的差异，以便他们各自向不同的人宣讲，不是不同的福音，而是（同一福音），\Person{伯多禄}向受割礼的人，\Person{保禄}向外邦人。那么，既然\Person{伯多禄}受责斥是因为他与外邦人同住之后，因顾及人的情面而主动离开他们的团体，这过错无疑是生活行为上的，而不是宣讲方面的。因为由此并不能看出他宣布了与造物主不同的神，或与（\Person{玛利亚}之子）不同的基督，或与复活不同的希望。\par

% structure: p0050 source=h2 target=section reason=relative-heading-rank; base=h2

\section{第二十四章 \Person{圣伯多禄}的进一步辩护。\Person{圣保禄}在教导上并不优于\Person{圣伯多禄}。第三层天上所启示的未能使他增补信德。异端者夸耀似乎蒙受了启示给\Person{圣保禄}的一些秘密。}

我并未有幸，或更确切地说，我并未承担那令人不快的任务，要使宗徒们彼此争吵。但是，既然我们那些极其悖逆的挑剔者刻意提出那项斥责，目的是使先前的道理受到怀疑，我将为\Person{伯多禄}辩护，大意是：甚至\Person{保禄}也说，他\Quote{成了对一切人成为一切——对犹太人成为犹太人}，对非犹太人成为非犹太人，\Quote{为要赢得一切人}。因此，他们是按时间、人和原因来指责某些做法，而这些做法他们自己在类似的时间、人和原因下也会毫不犹豫地去做。（\Emph{例如}）好像\Person{伯多禄}也曾指责\Person{保禄}，因为\Person{保禄}在禁止割损的同时，自己却给\Person{弟茂德}行了割损。不必理会那些审判宗徒的人！一个令人欣慰的事实是：\Person{伯多禄}在殉道的荣耀上与\Person{保禄}处于同一水平。现在，虽然\Person{保禄}曾被提升到第三层天，并被提到乐园里\BibleRef{格后 12:4}，在那里听见了某些启示，但这些启示根本不可能使他具备（教导）另一种道理的资格，因为它们的性质是那样，以至于不能传达给任何人。然而，如果那个不可言传的奥秘竟泄漏出来，被人知晓，并且如果有任何异端声称它本身也遵循这同一奥秘，那么，要么\Person{保禄}必须被指控泄露了秘密，要么必须证明后来另有其人曾被\Quote{提到乐园}，并获准公开说出\Person{保禄}连低声私语也不得说出的事。\par

% structure: p0052 source=h2 target=section reason=relative-heading-rank; base=h2

\section{第二十五章 宗徒们没有隐匿\Person{基督}托付给他们的任何道理宝库。\Person{圣保禄}公开将其全部道理交付给\Person{弟茂德}。}

但是，这里我们说过同样疯狂的是：他们固然承认宗徒们一无所不知，且所宣讲的道理彼此并不矛盾，但同时却坚持说，宗徒们并没有向所有人启示一切，因为他们向所有人公开宣讲了一些，而另一些则秘密地只向少数人揭示，这是因为\Person{保禄}甚至对\Person{弟茂德}说过这样的话：\Quote{啊！\Person{弟茂德}，要保管好所托付给你的；}\BibleRef{弟前 6:20}又说：\Quote{要持守那美好的托付。}\BibleRef{弟后 1:14}这托付是什么？难道它如此秘密，以至于被认为是一种新道理的特征？或者是那命令的一部分，他说：\Quote{我把这命令托付给你，我儿\Person{弟茂德}？}\BibleRef{弟前 1:18}以及那诫命的一部分，他说：\Quote{我在天主面前，就是那使万物生存的天主面前，并在\Person{耶稣基督}面前，就是那在本丢\Person{比拉多}前宣认过美好信仰的那位面前，嘱咐你：要持守这诫命？}\BibleRef{弟前 6:13}那么，什么是（这）诫命，什么是（这）命令？从上下文看来，显而易见，这句话并非对某种深奥的道理作隐晦的暗示，而是警告不要接受任何其他道理，除了那\Emph{弟茂德}从他本人那里听见的，就如我所理解的，公开地：\Quote{在许多见证人面前}是他的措辞。\BibleRef{弟后 2:2}现在，如果他们拒绝承认教会就是这些\Quote{许多见证人}所指的，那也无妨，既然\Quote{在许多见证人面前}产生的事不可能是秘密的。再者，\Person{保禄}希望他\Quote{把这些事托付给忠信的人，他们也能教导别人}\BibleRef{弟后 2:2}，这也不能被解释为存在一种隐秘福音的证明。因为当他说\Quote{这些事}时，是指他此刻所写的事。然而，对于隐秘的事，他会对与他共同知晓的人称它们为\Emph{那些事}，而不是\Emph{这些事}，因为那些事是隐而未现的。\par

% structure: p0054 source=h2 target=section reason=relative-heading-rank; base=h2

\section{第二十六章 宗徒们在所有情况下都将全部真理教导整个教会。毫无保留，也没有向偏爱的朋友作部分传达。}

此外，必然的后果是：对于那接受福音职务的人，主还会加上命令，不可在所有地方、不偏待人的方式施行这职务，这符合主的话：\BibleQuote{不要把珍珠丢在猪前，也不要把圣物给狗。}\BibleRef{玛 7:6} 主曾公开说话，\BibleRef{若 18:20}，没有暗示任何隐藏的奥迹。祂曾亲自命令，凡他们在黑暗中和秘密中所听到的，应在光明中并在房顶上宣讲出来。\BibleRef{玛 10:27} 祂也曾藉比喻预先表明，他们不可将仅一锭银子——即祂的一句话——秘密收藏而不生利息。\BibleRef{路 19:20} 祂又告诉他们，蜡烛通常不会放在斗底下，而是放在灯台上，为照亮屋里所有的人。\BibleRef{玛 5:15} 宗徒们若非忽略了这些事，就是未能明白，若他们不遵行，即隐藏了任何一部分光明——也就是天主的言语和基督的奥迹。我十分确信，他们不惧怕任何人——无论是犹太人还是外邦人的暴力；那么，他们既然在会堂和公共场所并未缄口不言，在教会中宣讲时必然更加自由。实际上，若不\BibleQuote{按次序叙述}\BibleRef{路 1:1}他们所要人相信的事，就不可能使犹太人皈依或使外邦人归信。更何况当教会在信德上已有长进时，他们更不会从中撤回任何东西，以便单独交托给少数其他人。虽然，可以这样说，即便在亲密朋友之间他们进行过某些讨论，但难以相信这些讨论竟会引入另一种与他们在各天主教会中所宣讲的不同的、相反的信仰规则——仿佛他们在教会中谈论一个天主，在家中谈论另一个；公开描述基督的一个本体，秘密描述另一个；向所有人宣布一种复活的希望，向少数人宣布另一种；虽然他们自己在书信中恳求众人都说一样的话，在教会中不可有分裂和纷争，\BibleRef{格前 1:10} 因为无论是保禄还是其他人，都宣讲相同的事。此外，他们记住了这句话：\BibleQuote{你们的话当是，是；非，非；除此之外，皆是出于邪恶。}\BibleRef{玛 5:37} 所以他们不应以不同方式处理福音。\par

% structure: p0056 source=h2 target=section reason=relative-heading-rank; base=h2

\section{第27章。即便宗徒们传授了全部真理，教会岂会在传递时有所不忠？这简直不可思议。}

因此，既然宗徒们不可能不知道他们所要宣告的全部信息，也没有向所有人宣告完整的信仰规则，那么让我们看看，也许当宗徒们简全地宣讲时，教会因自己的过失而传讲得与宗徒们不同。异端者们提出的所有不信任的理由，你都可以找到。他们记得教会如何被宗徒责备：\BibleQuote{无知的迦拉达人，谁迷惑了你们？}\BibleRef{迦 3:1} 以及 \BibleQuote{你们跑得好，谁阻挡了你们？}\BibleRef{迦 5:7} 还有书信开头：\BibleQuote{我希奇你们这么快就离开那以恩宠召你们的天主，去归从别的福音。}\BibleRef{迦 1:6} 他们还同样记得写给格林多人的信，说他们\BibleQuote{仍是属血肉的}，需要\BibleQuote{吃奶}，因为还\BibleQuote{不能吃干粮}；他们也\BibleQuote{自以为知道什么，其实按他们应当知道的，他们还不知道。}\BibleRef{格前 8:2} 当他们提出教会受责备的反对理由时，让他们设想教会也受了纠正；也让他们记住那些教会，宗徒为它们的信德、知识和行为\Quote{欢喜并感谢天主}，而这些教会直到今日仍与那些受责备的教会共同享有同一体制的特惠。\par

% structure: p0058 source=h2 target=section reason=relative-heading-rank; base=h2

\section{第28章。各处教会实质上相同的唯一信仰传承，证明传递基本上是真实而诚实的。}

那么，假定所有人都错了；宗徒作证错了；圣神并没有特别关照任何（教会）引导它进入真理，尽管基督差遣圣神正是为此目的，\BibleRef{若 14:26} 并且为此向圣父祈求，使祂成为真理的导师；\BibleRef{若 15:26} 再假定祂，天主的管家、基督的代理，忽略了祂的职务，允许教会在一段时间内对祂自己藉宗徒所宣讲的有不同的理解和信仰——难道如此众多、如此伟大的教会竟可能都错入了同一个信仰？偶然事件分布在许多人中不会产生同一结果。教义的错误必然在不同教会中产生不同的结果。然而，当在许多人中所存放的发现是同一的，那就不是错误的结果，而是传统的结果。那么，谁还敢轻率地说那传递传统的人错了呢？\par

% structure: p0060 source=h2 target=section reason=relative-heading-rank; base=h2

\section{第29章。真理并不依赖异端者的关照；真理在他们出现之前已自由传播。教会教义的优先性是其真理的标志。}

无论错谬以何种方式出现，它当然只在没有异端存在时才得以统治？真理竟要等待某些马尔基翁派和瓦伦提尼安派来释放它。在此期间，福音被错误地宣讲；人们错误地相信；如此众多的人被错误地受洗；如此多的信德工程被错误地施行；如此多的奇恩异能、如此多的神恩被错误地运作；如此多的司铎职务、如此多的服务被错误地执行；而一言以蔽之，如此多的殉道者错误地领受了他们的冠冕！若非错误且徒劳，那么天主的诸事在尚未知晓属于哪一位天主之前便已进行，这又如何可能呢？在基督被发现之前就有了基督徒？在真道之前就有了异端？并非如此；因为在一切情形下，真理先于其摹本，相似后于实在。然而，认为异端先于其自身的先在之道理，也够荒谬的——正是因为那道理本身预言了将有异端，而人必须防备它们！拥有这道理的教会曾收到书信——是的，那道理本身致信其自己的教会——\Quote{即使从天上来的天使，若向你们宣讲与我们先前所宣讲的不同的福音，当受诅咒。}\par

% structure: p0062 source=h2 target=section reason=relative-heading-rank; base=h2

\section{第三十章　异端相对较晚：马尔基翁的异端；关于他的若干个人事实；阿佩莱斯的异端；此人的品性；斐鲁墨；瓦伦提努斯；尼基狄乌斯和赫尔摩革乃。}

那时马尔基翁在哪里？那个本都的船主、热心钻研斯多葛主义的人？那时瓦伦提努斯在哪里？那个柏拉图主义的门徒？因为显然这些人生活在不太久以前——主要是在安东尼统治时期——并且他们起初在罗马教会、在真福埃留特略主教任期内信奉天主教会的道理，直到因他们那不倦的好奇心（他们甚至以此感染了弟兄们）而被多次驱逐。马尔基翁确实带着他并入教会的二百塞斯特斯离开了，而在最终被处以永久绝罚后，他们将教义的毒素散布各处。之后，马尔基翁确实表示悔改，并同意向他提供的条件——若他使所有他曾训练至丧亡的人恢复与教会的共融，他将获得和好；然而他被死亡阻止了。的确必须有异端；\BibleRef{格前 11:19}但这必然性并不因此使异端成为好事。仿佛也必须有邪恶一样！甚至主也必须被出卖；但出卖者是有祸的！\BibleRef{谷 14:21}因此，没有人可从这一点为异端辩护。如果我们也必须谈及阿佩莱斯的传承，他远非\Quote{旧派人物}，不像他的导师和塑造者马尔基翁；他反而放弃了马尔基翁的节制，与一女子为伴，并退到亚历山大里亚，远离他最克己的主人。数年后从未处返回，毫无改善，只是不再是一个马尔基翁派，他又依恋另一个女子，即童女斐鲁墨（我们已提到过她），她后来竟成了一名大妓女。被她那活跃的灵所欺骗，他将从她那里学到的\Emph{启示}写了下来。现今还有人记得它们——他们自己的门徒和继承者——因此不能否认它们年代较晚。但其实，正如主所说，他们被自己的行为所定罪。\BibleRef{玛 7:16}因为既然马尔基翁将新约与旧约分开，他（必然）后于他所分开的，因为他只能分开先前（已）联合的。既然在分开之前已是联合的，后来分开的事实也证明了实行分开者之后出现。同样，瓦伦提努斯通过不同的解释和公认的修正，明确以先前的缺陷为由进行这些更改，因此证明了文献的差异。这些真理的败坏者，我们提及他们是因为他们比其他人更臭名昭著和更公开。然而，有一个名叫尼基狄乌斯的人，还有赫尔摩革乃，以及其他几个，仍在继续从事扭曲主之道的勾当。让他们告诉我，他们以何权威而来！若他们宣讲的是另一个天主，他们怎能使用他们所反对的那位天主的事物、著作和名字呢？若是同一位天主，为何以另一种方式对待他？让他们证明自己是新宗徒！让他们主张基督第二次降临，第二次亲自教导，被钉十字架两次，死两次，复活两次！因为宗徒正是这样描述（基督生命中的事件次序）的；也正因如此，主惯常制造他的宗徒——即赐给他们（能力）去行他自己所行的同样奇迹。因此，我也愿意他们的异能也显出来；只是我认为他们最大的异能便是他们与宗徒相敌的扭曲之行。因为当宗徒们曾使人从死里复活时，这些人却使活着的人死去。\par

% structure: p0064 source=h2 target=section reason=relative-heading-rank; base=h2

\section{第三十一章　真理在先，错谬在后，作为其歪曲；基督的比喻将好种子撒在无用的莠子之前。}

然而，请允许我从这一题外话转回来，讨论真理的优先性和虚假的相对晚近性，甚至从那比喻中为我的论点寻找支持：那比喻首先提到主播下麦子的好种子，随后才引入仇敌魔鬼用无用的野燕麦稗子来掺杂庄稼。因为这里形象地描述了教义的不同，在其他经文中天主的圣言也被比作种子。因此，从实际顺序来看，显而易见：首先传下来的属于主，是真理；而后来引入的则是陌生和虚假的。这一判断将屹立不倒，反对所有后来的异端，这些异端本身并不具有同源知识的连贯品质——无法声称真理在他们那边。\par

% structure: p0066 source=h2 target=section reason=relative-heading-rank; base=h2

\section{第三十二章：没有异端声称从宗徒继承。新教会仍是宗徒的，因为她们的信仰是宗徒所教导和传授的。挑战异端出示任何宗徒凭证。}

但如果有任何异端胆敢置身于宗徒时代之中，以便显得是由宗徒传下来的，因为它们在宗徒时代就已存在，我们可以说：让它们出示其教会的原始记录；让它们展开主教的名单，从开始就按正当顺序传承下来，以至于那第一位主教能够证明其祝圣者和前任是某位宗徒或宗徒弟子——而且是坚持与宗徒同在的人。因为宗徒教会就是这样传递她们的登记册的：如\Place{斯弥纳}教会记载\Person{波利卡普}是由\Person{若望}安置在那里的；又如\Place{罗马}教会记载\Person{克莱孟}是由\Person{伯多禄}以类似方式祝圣的。同样地，其他教会也展示她们的杰出人物，这些人被宗徒任命于他们的主教职位，被视为宗徒种子的传递者。让异端们也炮制同样的东西。因为在他们的亵渎之后，还有什么是不允许他们尝试的呢？但即使他们炮制成功，也前进不了半步。因为他们的教义，在与宗徒的教义比较之后，会因其自身的多样性和矛盾性而宣告，其作者既不是宗徒也不是宗徒弟子；因为，正如宗徒绝不会教导自相矛盾的事，宗徒弟子也不会教导与宗徒不同的道理，除非他们从宗徒接受教导后却去宣讲相反的道理。因此，这些教会将成为检验他们的证据：这些教会虽然不从宗徒或宗徒弟子那里得到其奠基者（因为它们成立得晚得多——事实上它们每天都在成立），但由于它们持守同一信仰，因教义相通而被视为同样具有宗徒性。那么，让所有异端在受到我们宗徒教会的这两项检验时，提出证据表明它们如何自认为属于宗徒。但事实上它们既不是，也无法证明它们所不是的。它们也不被那些与宗徒有某种联系的教会接纳进入和平的关系和共融，因为它们在信仰奥秘上的差异使它们本身毫无宗徒性质。\par

% structure: p0068 source=h2 target=section reason=relative-heading-rank; base=h2

\section{第三十三章：现在的异端（神圣作者所记载的稗子的幼苗）已在圣经中被定罪。后来的异端从早期的异端衍生，这在几个实例中得到了追溯。}

此外，我还要审视那些教义本身，它们既存在于宗徒时代，也被上述宗徒揭露和谴责。因为通过这种方法，它们将更容易被废弃——当发现它们甚至在当时就已存在，或者至少是当时就已存在的（莠子）的幼苗时。保禄在致格林多人前书中，指责那些否认和怀疑复活的人。\BibleRef{格前 15:12}这种观点是撒杜塞人的特有主张。然而，其一部分为玛尔其翁、阿培肋斯、瓦伦提诺以及所有其他攻击复活者所坚持。同样，在致迦拉达人书中，他痛斥那些遵守和辩护割损以及（梅瑟）法律的人。\BibleRef{迦 5:2}厄俾翁的异端正是如此。他也在训诲弟茂德时责备那些\Quote{禁止结婚}的人。\BibleRef{弟前 4:3}现在，这是玛尔其翁及其追随者阿培肋斯的教导。（宗徒）同样打击那些说\Quote{复活已成过去}的人。\BibleRef{弟后 2:3}瓦伦提诺派的人正是如此自称。当他再次提到\Quote{无尽的族谱}时\BibleRef{弟前 1:4}，我们认出瓦伦提诺，在他的体系中，某个永世（无论他是谁），有一个新名字，而且不止一个，由其恩宠产生\Quote{感知}和\Quote{真理}，然后这些又同样产生\Quote{言辞}和\Quote{生命}，这些后又生出\Quote{人}和\Quote{教会}。从这最初的八位，又生出另外十位永世，然后有十二位以奇妙的名称出现，完成这三十永世的纯粹故事。同一宗徒，在反对那些\Quote{受元素束缚}的人时\BibleRef{迦 4:9}，指出赫尔莫革内斯的某个教义，他引入无始的物质，然后将其与无始的天主相比较。通过这样使元素的母亲成为女神，他能够\Quote{受束缚}于一个他认为与天主同等的存在。然而，若望在默示录中奉命惩治那些\Quote{吃祭偶像之物}和\Quote{行淫乱}的人。\BibleRef{默 2:14}现在甚至还有另一种尼苛劳党人。他们的被称为加约异端。但在他的书信中，他特别指出那些\Quote{否认基督成了肉身}的人\BibleRef{若一 4:3}和那些拒绝相信耶稣是天主子的人为\Quote{假基督}。玛尔其翁坚持一个教义，厄俾翁另一个。然而，西满巫术的教义，即灌输崇拜天使，本身实际上被列为偶像崇拜，并被宗徒伯多禄在西满本人身上定罪。\par

% structure: p0070 source=h2 target=section reason=relative-heading-rank; base=h2

\section{第三十四章 关于天主造物主并无早期争论；起初未引入第二位神。异端被圣经的判决和沉默同样定罪。}

我想，这些都是各种虚假教义，它们（如宗徒们亲自告知我们的）存在于他们自己的时代。然而，在如此众多的真理曲解中，我们发现没有一个学派对作为万物造物主的天主提出争议。没有人胆敢猜测第二位神。人们更容易怀疑圣子而不是圣父，直到玛尔其翁在造物主之外引入另一位只属于善的神。阿培肋斯把造物主当作某个无名的荣耀天使，这天使属于那上等的神，即（按他的说法）法律和以色列的神，并声称他是火。瓦伦提诺散布他的永世，并将一个永世的罪追溯至天主造物主的产生。当然，除了这些，且在它们之前，没有人得到关于天主性体的真理的启示；他们从魔鬼那里获得这种特殊的荣誉和更丰厚的恩宠，我们毫不怀疑，因为魔鬼甚至在这方面也要与天主较量，好藉其教义的毒害，亲自做成主所说不可能做到的事——使\Quote{门徒高于师傅}。\BibleRef{路 6:40}因此，让所有异端的整体为它们自己选择出现的时间，只要这\Emph{何时}是不重要的；也承认它们不属于真理，并且（当然）那些在宗徒时代不存在的不可能与宗徒有任何关联。如果它们确实存在过，它们的名字也会留存下来，以便同样压制它们。那些确实存在于宗徒时代的（异端），单是提及它们就被定罪。那么，如果那些在宗徒时代尚处于粗糙形式的异端，现在被发现是相同的，只是形状更为精致，它们正是因这一情况而得定罪。或者，如果它们并不相同，而是以后以不同形式出现，并仅仅从它们那里接受了某些教义，那么，由于在教导上与它们一致，它们必定因上述关于晚近出现的定义而分受其定罪，这定义从一开始就面对我们。即使它们完全没有参与被定罪的教义，它们仅因时间理由就已受审判，因为它们甚至未被宗徒提及而更加虚假。因此，我们有更坚定的确据，即这些就是当时被预告将要出现的（异端）。\par

% structure: p0072 source=h2 target=section reason=relative-heading-rank; base=h2

\section{第三十五章 让异端者以明确可理解的证据维护其主张。这是解决他们问题的唯一方法。天主教徒始终诉诸可追溯至宗徒来源的证据。}

根据这些定义，我们已将所有异端驳倒。现在，让它们也大胆地提出与我们的教义相对立的类似规则吧，无论它们是晚于宗徒还是与宗徒同时代，只要与宗徒的教训不同；并且它们已受到宗徒们普遍的或具体的谴责。因为它们否认（我们教义的）真理，就应当证明它也是异端，可用与它们自身被驳倒相同的规则来反驳；同时向我们指明，我们必须在哪里寻求真理——现在显然在它们那里并不存在真理。我们的体系在年代上不晚于任何体系；相反，它比所有体系都更早；这一事实将证明那真理的出处——它处处占据首位。再者，宗徒们从未谴责它；他们反而捍卫它——这一事实表明它出自他们本身。因为那教义，当他们谴责了所有异己见解时，却从不加以谴责，这就表明它是他们的，并且他们也在此基础上捍卫它。\par

% structure: p0074 source=h2 target=section reason=relative-heading-rank; base=h2

\section{第三十六章——宗徒的教会即宗徒的声音。让异端者检验各教会的宗徒传承主张，每案皆无可争议。罗马教会具有双重宗徒传承；其早期卓越与优秀。异端，因其歪曲真理，与此相连。}

来吧，你们这些渴望更佳探求的人，如果你们愿意将此心用于你们的救恩事业，就去遍历宗徒的教会吧，在那裡，宗徒们的宝座至今仍以其陈设突显出来，在那裡，他们自己亲笔的著作仍被诵读，发出每个人的声音，代表每个人的面容。阿哈雅离你们很近，（在那裡）你们能找到格林多。既然你们离马其顿不远，你们有斐理伯；（还有）你们有得撒洛尼。既然你们能渡海到亚细亚，你们能到厄弗所。既然，此外，你们接近意大利，你们有罗马——从罗马甚至传到我们手中（宗徒们的）权威。她的教会是多么幸福，宗徒们把他们的全部教义连同他们的血一起倾注其中！在那里，伯多禄经历与主相似的苦难！在那里，保禄赢得像若翰一样的死亡中的冠冕——在那里，若望宗徒先被投入沸腾的油中而无损伤，然后被遣往他的岛屿流放！看看她学到了什么，教导了什么，甚至与（我们的）非洲教会有什么共融！她承认一位主天主，宇宙的创造者，以及耶稣基督（生于）童贞玛利亚，是创造者天主的圣子；并承认肉身的复活；她把法律与先知同福音作者和宗徒的著作合为一卷，从中汲取信仰。她用（圣洗的）水封印这信仰，以圣神装备，以圣体喂养，以殉道鼓励，并且对于这样（遵守的）一种纪律，她不容许任何反对者。这就是那纪律——我不再说它预言了异端的到来，而是说异端从它产生。然而，它们不属于她，因为它们与她相对立。甚至粗糙的野橄榄也从丰硕、饱满、纯正的橄榄的胚芽中长出；也从最甜美、最甘甜的无花果的种子中生出空洞无用的野无花果。同样，异端也来自我们的植物，尽管不属于我们的种类；它们来自真理的谷粒，但由于它们的虚假，只显示野生的叶子。\par

% structure: p0076 source=h2 target=section reason=relative-heading-rank; base=h2

\section{第三十七章——异端者并非基督徒，而是基督教训的歪曲者，无权要求基督信仰的圣经。圣经是托付给教会并受教会谨慎保管的寄存物。}

既然情况如此，为了让真理被判定属于我们——\Quote{凡按这准则而行的人}，这准则就是教会从宗徒、宗徒从基督、\Emph{以及}基督从天主所传授的——我们的立场就清楚了，即决定异端者不应被允许使用圣经，因为我们在没有圣经的情况下就能证明他们与圣经无关。因为他们既是异端者，就不可能是真正的基督徒，因为他们所追求的东西不是来自基督，而是出于自己的私意选择，并且由于这种追求而得名为异端。这样，他们既然不是基督徒，就没有获得对基督徒圣经的权利；因此完全可以对他们说：\Quote{你们是谁？你们何时何地来的？既然你们不是属我的，你们与属我的有什么关系？确实，马西翁，你凭什么权利砍我的树木？瓦伦提诺，你得到谁的许可改变我泉源的流向？阿佩莱斯，你凭什么能力移动我的地界？这是我的产业。你们其余的人，为什么在这里任意播种和放牧？这（我说）是我的产业。我早已拥有它；我在你们之前就拥有它。我从原主自己手里持有可靠的产权证书，这产业属于他们。我是宗徒的继承人。正如他们精心制定他们的遗嘱并将其托付给信托人，并恳求（信托人忠实履行职责），我同样持守它。至于你们，他们确实一直视你们为被剥夺继承权的人，将你们作为外人——作为敌人加以拒绝。然而，如果异端者不是由于他们与宗徒教义的差异——即每个人凭自己的私意或提出或接受反对宗徒的教义——他们又何以是宗徒的外人和敌人呢？}\par

% structure: p0078 source=h2 target=section reason=relative-heading-rank; base=h2

\section{第三十八章——教会与圣经的和谐。异端者已篡改圣经，并加以肢解和改动。公教徒从不改变圣经，圣经始终为他们作证。}

凡发现教义分歧之处，就应视为圣经及其解释的败坏必然存在于此。那些意欲教导不同教义的人，必须承担以不同方式安排教导工具的责任。他们不可能以其他方式实现其教导的多样性，除非其所用以教导的手段也存在差异。正如在他们那里，教义的败坏若无其工具的败坏便不可能成功，同样，在我们这里，教义的纯正若无管理教义之工具的纯正亦不可能获得。那么，在我们的圣经中有什么与我们对立呢？我们引入了什么属于我们自己的东西，以至于我们必须将其移除，或增添，或更改，以便恢复其中与我们对立的任何东西的自然纯正，而这些对立之物又载于圣经之中？我们自己是什么样，圣经从起初就是（并且一直是）什么样。我们由圣经而有存在，早于任何其他方式，早于它们被你们篡改。既然所有篡改必被视为后来的过程，明确的原因在于它产生于竞争，而竞争在任何情况下都不会先于它所模仿的对象，也不是与它所模仿的对象同生的，因此，对于每个有理智的人来说，我们，自始至终存在并且是最初的，似乎已在圣经中引入了任何败坏的文本，这是不可信的，正如那些在时间上较晚且与圣经对立的人实际上并没有引入败坏一样。一个人亲手曲解圣经，另一个人则以其解释曲解圣经的意思。因为虽然\Person{瓦伦廷}似乎使用了整部圣经，但他比\Person{马西昂}以更狡猾的心思和技巧对真理施以暴力。\Person{马西昂}明确而公开地使用刀，而非笔，因为他按自己的主题删改了圣经。然而，\Person{瓦伦廷}却避免了这样的删改，因为他不是为使圣经符合自己的主题而虚构圣经，而是使自己的材料适应圣经；然而他却取走了更多，添加了更多，通过去除每个特定词语的适当意义，并添加了并无真实存在的虚构安排。\par

% structure: p0080 source=h2 target=section reason=relative-heading-rank; base=h2

\section{第三十九章 圣保禄所称的属灵的邪恶，由外邦作者和异端以并非不同的方式展示。圣经尤其易受异端操纵。为异端提供材料，正如\Person{维吉尔}已成为与原文旨趣不同的文学剽窃的底本。}

这些正是\Quote{属灵的邪恶}的巧妙技艺，我们，我的弟兄们，也理当期望与之\Quote{搏斗}，作为信德所需，为使被选者显明，被弃者得以发现。因此，它们具有影响力，并且在思考和编造错谬方面有技巧，这不应当被惊奇，仿佛这是一个困难而无法解释的过程，因为在世俗著作中也随手可得到类似技巧的例子。你们看到在我们今天，从\Person{维吉尔}的作品中编出的一个完全不同的故事，主题按诗句编排，诗句也按主题编排。简而言之，\Person{霍西狄乌斯·革塔}完全从\Person{维吉尔}的作品中剽窃了他的悲剧\WorkTitle{美狄亚}。我的一位近亲在其闲暇之作中，从同一诗人创作了\WorkTitle{塞贝斯的画板}。基于同样的原理，那些\Emph{劣诗作者}通常被称为\Emph{荷马拼缀诗人}，即\Quote{荷马零散片段的收集者}，他们以拼缝的方式，从荷马的诗行中将自己的作品缀合成一体，由许多从各处（杂乱地）拾取的片段拼凑而成。现在，无疑，神圣的圣经在这种技巧方面拥有更丰富的各种资源。我说圣经甚至是由天主的旨意如此安排，以便为异端提供材料，这并不冒自相矛盾之险，因为我读到\BibleQuote{必须有异端}\BibleRef{格前 11:19}，而没有圣经则不能有异端。\par

% structure: p0082 source=h2 target=section reason=relative-heading-rank; base=h2

\section{第四十章 偶像崇拜的精神与异端的精神并无区别。在偶像崇拜的礼仪中，撒殚模仿并歪曲了旧约圣经的神圣制度。基督宗教的圣经被他通过各类异端的曲解所败坏。}

将出现一个问题：谁应当解释那些有利于异端的经文的含义？当然，是魔鬼，他拥有那些歪曲真理的诡计，并且通过他偶像的神秘仪式，甚至与天主圣事的基本部分竞争。他同样给人施洗——即他自己的信徒和忠诚的追随者；他应许藉着（自己的）洗涤池除掉罪过；若我记忆无误，那里的密特拉（在撒殚的王国中）在他士兵的前额上做记号；也举行献饼礼，并引入一个复活的形象，并在刀剑前编织一个冠冕。我们又当如何说撒殚将其大司祭限于只娶一次婚？他也有自己的贞女；他也有自己的精进修德者。假设现在我们心中思索努玛·庞皮里乌斯的迷信，并思考他的司祭职务、徽章和特权，还有他的祭献服务，以及祭献本身的器具和器皿，以及他赎罪和誓愿的奇特礼仪：难道我们不清楚魔鬼模仿了犹太法律那众所周知的严厉吗？因此，既然他在表达其偶像崇拜事务中所表现的争胜心，正是基督圣事管理所包含的那些事，那么，理所当然，同一位存有者，仍然拥有同样的才智，既专心致力于，也成功地调整适应于他那亵渎且敌对的信条——神性事物和基督徒圣人们的文献：他的解释取自他们的解释，他的言词取自他们的言词，他的比喻取自他们的比喻。因此，任何人都不应怀疑：或是由魔鬼引入了\Quote{属灵的邪恶，}异端也由此而来；或是异端与偶像崇拜之间有任何真正区别，因为它们与偶像崇拜同属同一作者和同一工作。他们要么假装有另一位神与造物主对立，要么即使承认造物主是唯一的天主，他们也将祂当作与祂真实面貌不同的存有者。结果是，他们所说关于天主的每一谎言，在某种意义上都是一种偶像崇拜。\par

% structure: p0084 source=h2 target=section reason=relative-heading-rank; base=h2

\section{第四十一章 异端者的行为：其轻浮、世俗与失序。其妇女臭名昭著的放荡。}

我不可略过关于异端者行为的描述——多么轻浮、多么世俗、多么人性，没有严肃、没有权威、没有纪律，正如适合他们的信条。首先，谁是慕道者、谁是信徒尚存疑问；他们所有人进出相同、听闻相同、祈祷相同——甚至外邦人，若有任何这等偶入其间者亦然。\Quote{他们将圣物丢给狗，并将他们的珍珠，}虽然（当然）并非真的，\Quote{他们将抛给猪。}简朴性必须在于推翻纪律，而我们对此的关注，他们却称之为娼妓行为。平安他们也随意与任何人草率成就；因为对他们来说，无论他们的主题处理如何不同，只要他们能共同密谋攻取独一真理的堡垒，就无关紧要。所有人都自负，所有人都提供知识。他们的慕道者在尚未完全受教之前就已完美。这些异端者的妇女，多么放荡！因她们竟敢教导、辩论、施行驱魔、从事治病——甚至可能施洗。他们的圣秩授予，漫不经心、反复无常、变动不定。有时他们让\Emph{新手}任职；有时让受世俗职业束缚的人；有时让从我们这里背教的人，为的是以虚荣束缚他们，因他们无法以真理。在叛党营中，晋升无处更易，在那里，仅仅身在营中就是首要的服事。于是发生这样的事：今天一人是他们的主教，明日又是另一人；今天是执事，明日成了读经员；今天是司铎，明日成了平信徒。因为他们甚至将司铎的职务强加于平信徒。\par

% structure: p0086 source=h2 target=section reason=relative-heading-rank; base=h2

\section{第四十二章 异端者的工作是拆毁和摧毁，而非建立和提升。异端者甚至不固守自己的传统，反而对自身创始人也心怀异见。}

然而，论及圣言的职务，我该说什么呢？因为他们的事业不是归化外邦人，而是颠覆我们的人。这倒是他们所追求的荣耀：设计使站立者跌倒，而非扶起跌倒者。因此，既然他们为自己设定的工作不是来自建立自己的团体，而是来自拆毁真理，他们就破坏我们的建筑，为要建立自己的。只要夺走他们的梅瑟法律、先知和造物主的神性，他们就没有别的反对之辞可说了。结果是，他们更易成就对站立房屋的毁坏，而非对已倒塌废墟的建造。他们只有在瞄准这些目标时，才显出谦卑、温和和恭敬。否则他们对自身领袖也不知尊敬。因此，人们以为异端者中很少发生分裂，因为即使存在分裂，也不明显。然而，他们的合一本身即是分裂。若他们不在自己中间也偏离自身的规条，我就大错特错了，因为每个人，正因符合自己的性情，修改他所接受的传承，其方式与传授者当初按自己意愿塑造它们时相同。事情的发展同时承认了其性质及其产生的方式。华伦提努斯所允许的，对华伦提努派也是允许的；马尔基昂所做的，对马尔基昂派也是公平的——甚至按自己的喜乐创新信仰。总之，所有异端，若仔细查究，都被发现在许多细节上甚至与自己的创始人怀抱异议。他们大多数甚至没有教会。无母、无家、无信经，被遗弃，他们在其本质的卑贱中四处游荡。\par

% structure: p0088 source=h2 target=section reason=relative-heading-rank; base=h2

\section{第四十三章 异端者偏爱松散团体。他们的教导产生不虔敬，这与天主教的真理截然相反，后者在宗教礼仪和实际生活中增进对天主的敬畏。}

也有人注意到，异端者与术士、江湖骗子、占星家和哲学家交往极其频繁；原因是他们都是致力于探索奇闻异事的人。\Quote{你们求，必要给你们}，这句话处处萦绕在他们心头。因此，从他们行为的本质，就可以估量他们信德的质量。在他们的纪律中，我们看到了他们教义的索引。他们说天主是不可畏惧的；因此在他们的眼中，一切事物都是自由和不受约束的。但是，除了天主不\Emph{存在}的地方，哪里才会不畏惧天主呢？哪里没有天主，哪里也就没有真理。没有真理的地方，自然也就会有他们那样的纪律。但是，哪里有天主，哪里就存在\Quote{敬畏天主是智慧的开始}。哪里有敬畏天主，哪里就有严肃、可敬而又深思熟虑的勤勉，以及焦虑的谨慎和周全的接纳（进入圣职），有安全守护的共融，有良好服务后的晋升，有严格的服从（权威），有虔诚的参与，有谦逊的举止，有合一的教会，以及天主\Emph{在}万物之中。\par

% structure: p0090 source=h2 target=section reason=relative-heading-rank; base=h2

\section{第四十四章 异端降低对基督的尊重，并摧毁对其伟大审判的一切敬畏。异端教导对这信仰庄严信条的影响。本论著是我们作者其他某些反异端著作的引言。}

这些存在于我们中间的更严格纪律的证据，是真理的额外证明。任何记念那将来审判的人，都不能安全地偏离它，那时\BibleQuote{我们众人都应出现在基督的审判台前}\BibleRef{格后 5:10}，以首先为我们自己的信德交出账目。那么，那些玷污了那信德——即基督托付给他们的童贞——因与异端者行奸淫的人，将说什么呢？我想他们会声称，主和宗徒从未向他们下达过任何针对败坏和乖谬教义攻击他们的禁令，或关于避免和憎恶这些教义的命令。或许（主和宗徒）会承认责任在于他们自己和他们的门徒，因为没有事先警告和教导我们！此外，他们还会大肆吹嘘每个异端导师的崇高权威——他们如何有力地增强对其自身教义的信仰；如何复活死人、治愈病人、预言未来，从而使他们配得被视为宗徒。好像这警告没有写在圣经中：将有许多人来，甚至能行最大的奇迹，以维护他们腐败宣讲的骗局。因此，他们配得被宽恕！然而，如果有人记念主和宗徒的著作和谴责，在信德的完整中站稳，我想\Emph{他们}将冒失去宽恕的巨大风险，当主回答说：\Quote{我明明预先警告过你们，会有假师傅以我的名义，以及以先知和宗徒的名义来教导；我也曾命令我的门徒，要他们向你们宣讲同样的事。至于你们，当然不能指望你们会相信我！我当初将福音和上述规诫（生活和信德的）交给我的宗徒；但后来我很乐意在其中做出相当大的改变！我曾应许复活，甚至肉身的复活；但再想想，我意识到我可能无法履行我的应许！我曾显示自己是由童贞女所生；但这后来在我看来是一件可耻的事。我曾说过那创造太阳和雨水的圣父是我的父亲；但另一位更好的父亲收养了我！我曾禁止你们听从异端者；但在这一点上我错了！}这样的（亵渎）思想，很可能进入那些偏离正路、不保护真信德免受危险之人的头脑。的确，在目前的场合，我们的论著采取了反异端的一般立场，表明它们都必须根据明确、公平和必要的规则而被驳斥，无需与圣经比较。其余的事情，若天主恩准，我们将在单独的论著中准备回答某些异端。愿那些因相信真理而抽空阅读这些篇章的人，得享平安，以及我们天主耶稣基督的恩宠，直到永远。\par

\SourceRecord{Translated by Peter Holmes.；Ante-Nicene Fathers；Vol. 3.；Edited by Alexander Roberts, James Donaldson, and A. Cleveland Coxe.；Buffalo, NY: Christian Literature Publishing Co.,；1885.；<http://www.newadvent.org/fathers/0311.htm>.}
